% ================================================================================
% MANUAL DE INSTALACION — GAMMA
% Puesta en marcha, actualizacion y operacion de los contenedores.
% Sustituye a la pestana "Despliegue" que existia en la seccion Referencia.
%
% Compilar con:  latexmk -pdf manual_instalacion.tex
% ================================================================================
% ================================================================================
% PREAMBULO COMPARTIDO — Manuales de GAMMA
%
% Lo usan los tres manuales (usuario, tecnico e instalacion) para que compartan
% tipografia, paleta y recuadros. Se incluye con % ================================================================================
% PREAMBULO COMPARTIDO — Manuales de GAMMA
%
% Lo usan los tres manuales (usuario, tecnico e instalacion) para que compartan
% tipografia, paleta y recuadros. Se incluye con % ================================================================================
% PREAMBULO COMPARTIDO — Manuales de GAMMA
%
% Lo usan los tres manuales (usuario, tecnico e instalacion) para que compartan
% tipografia, paleta y recuadros. Se incluye con \input{../preambulo} desde la
% carpeta de cada manual.
%
% La paleta sale de las variables de frontend/src/style.css: la interfaz es
% enteramente acromatica, y los manuales tampoco introducen color.
% ================================================================================
\documentclass[11pt, letterpaper]{article}

\usepackage[spanish,mexico]{babel}
\usepackage[utf8]{inputenc}
\usepackage[T1]{fontenc}

% Tipografia: Inter, la mas cercana en TeX Live a la Geist que usa la
% aplicacion. Sin serifas en todo el documento, igual que la interfaz.
\usepackage[sfdefault]{inter}
\usepackage[scaled=0.86]{beramono}

\usepackage[top=2.4cm, bottom=2.4cm, left=2.6cm, right=2.6cm]{geometry}
\usepackage{graphicx}
\usepackage{xcolor}
\usepackage{enumitem}
\usepackage{fancyhdr}
\usepackage{titlesec}
\usepackage{tcolorbox}
% `enhanced` y `borderline west` viven en skins; `breakable` en su propia
% libreria. Sin declararlas, tcolorbox ignora esas claves en silencio y los
% recuadros pierden la barra lateral que los distingue.
\tcbuselibrary{skins,breakable}
\usepackage{booktabs}
\usepackage[hidelinks]{hyperref}

% ── Paleta ──────────────────────────────────────────────────────────────────
% Tomada de las variables de src/style.css. La interfaz es enteramente
% acromatica —todos sus colores son oklch(L 0 0)— asi que el manual tampoco
% introduce color: solo grises, con el texto como unico elemento de contraste.
\definecolor{fg}{HTML}{0A0A0A}        % --foreground
\definecolor{fgsoft}{HTML}{171717}    % --primary
\definecolor{muted}{HTML}{737373}     % --muted-foreground
\definecolor{surface}{HTML}{F5F5F5}   % --muted / --secondary
\definecolor{surfacealt}{HTML}{FAFAFA}% --sidebar
\definecolor{line}{HTML}{E5E5E5}      % --border

\titleformat{\section}{\LARGE\bfseries\color{fg}}{\thesection}{0.7em}{}
\titlespacing*{\section}{0pt}{1.6em}{0.7em}
\titleformat{\subsection}{\large\bfseries\color{fgsoft}}{\thesubsection}{0.6em}{}
\titlespacing*{\subsection}{0pt}{1.3em}{0.4em}
\titleformat{\subsubsection}[runin]{\bfseries\color{fgsoft}}{}{0pt}{}[.]

\pagestyle{fancy}
\fancyhf{}
\lhead{\footnotesize\color{muted}Manual de usuario}
\rhead{\footnotesize\color{muted}\thepage}
\renewcommand{\headrulewidth}{0.4pt}
\renewcommand{\headrule}{\color{line}\hrule height 0.4pt}

\setlength{\parindent}{0pt}
\setlength{\parskip}{0.62em}

% Enlaces y referencias en el mismo gris del cuerpo: sin subrayados ni azules.
\hypersetup{colorlinks=false}

% ── Recuadros ───────────────────────────────────────────────────────────────
% Se distinguen por una barra lateral y el peso del titulo, no por color.
\newtcolorbox{nota}[1][Nota]{
  enhanced, breakable, colback=surfacealt, colframe=surfacealt,
  borderline west={2pt}{0pt}{line},
  coltitle=fg, fonttitle=\bfseries\footnotesize\sffamily,
  title=#1, boxrule=0pt, arc=0pt, outer arc=0pt,
  left=10pt, right=10pt, top=7pt, bottom=7pt, toptitle=1pt, bottomtitle=3pt}

\newtcolorbox{aviso}[1][Importante]{
  enhanced, breakable, colback=surface, colframe=surface,
  borderline west={2pt}{0pt}{fg},
  coltitle=fg, fonttitle=\bfseries\footnotesize\sffamily,
  title=#1, boxrule=0pt, arc=0pt, outer arc=0pt,
  left=10pt, right=10pt, top=7pt, bottom=7pt, toptitle=1pt, bottomtitle=3pt}

% ── Captura de pantalla con marcador de posicion ────────────────────────────
% Uso: \captura{nombre-archivo}{Pie de la figura}
% Si manual/capturas/<nombre>.png existe, la inserta. Si no, dibuja el hueco.
\newcommand{\captura}[2]{%
  \begin{figure}[htbp]
  \centering
  \IfFileExists{capturas/#1.png}{%
    % Ancho completo, pero acotado en alto: una captura de ventana completa
    % puede ser muy alta y desbordaria la pagina. keepaspectratio evita que se
    % deforme cuando manda la restriccion de altura.
    {\color{line}\fboxrule=0.6pt\fboxsep=0pt\fbox{%
      \includegraphics[width=0.94\textwidth, height=0.62\textheight, keepaspectratio]{capturas/#1.png}}}%
  }{%
    \fcolorbox{line}{surfacealt}{%
      \begin{minipage}[c][3.1cm][c]{0.93\textwidth}
        \centering
        {\footnotesize\bfseries\color{muted}CAPTURA PENDIENTE}\\[5pt]
        {\ttfamily\footnotesize\color{fgsoft} capturas/#1.png}\\[4pt]
        {\footnotesize\color{muted} #2}
      \end{minipage}}%
  }
  \caption{#2}
  \label{fig:#1}
  \end{figure}%
}

% Pies de figura discretos
\usepackage{caption}
\captionsetup{font={small,color=muted}, labelfont={bf,color=muted}, skip=6pt}

 desde la
% carpeta de cada manual.
%
% La paleta sale de las variables de frontend/src/style.css: la interfaz es
% enteramente acromatica, y los manuales tampoco introducen color.
% ================================================================================
\documentclass[11pt, letterpaper]{article}

\usepackage[spanish,mexico]{babel}
\usepackage[utf8]{inputenc}
\usepackage[T1]{fontenc}

% Tipografia: Inter, la mas cercana en TeX Live a la Geist que usa la
% aplicacion. Sin serifas en todo el documento, igual que la interfaz.
\usepackage[sfdefault]{inter}
\usepackage[scaled=0.86]{beramono}

\usepackage[top=2.4cm, bottom=2.4cm, left=2.6cm, right=2.6cm]{geometry}
\usepackage{graphicx}
\usepackage{xcolor}
\usepackage{enumitem}
\usepackage{fancyhdr}
\usepackage{titlesec}
\usepackage{tcolorbox}
% `enhanced` y `borderline west` viven en skins; `breakable` en su propia
% libreria. Sin declararlas, tcolorbox ignora esas claves en silencio y los
% recuadros pierden la barra lateral que los distingue.
\tcbuselibrary{skins,breakable}
\usepackage{booktabs}
\usepackage[hidelinks]{hyperref}

% ── Paleta ──────────────────────────────────────────────────────────────────
% Tomada de las variables de src/style.css. La interfaz es enteramente
% acromatica —todos sus colores son oklch(L 0 0)— asi que el manual tampoco
% introduce color: solo grises, con el texto como unico elemento de contraste.
\definecolor{fg}{HTML}{0A0A0A}        % --foreground
\definecolor{fgsoft}{HTML}{171717}    % --primary
\definecolor{muted}{HTML}{737373}     % --muted-foreground
\definecolor{surface}{HTML}{F5F5F5}   % --muted / --secondary
\definecolor{surfacealt}{HTML}{FAFAFA}% --sidebar
\definecolor{line}{HTML}{E5E5E5}      % --border

\titleformat{\section}{\LARGE\bfseries\color{fg}}{\thesection}{0.7em}{}
\titlespacing*{\section}{0pt}{1.6em}{0.7em}
\titleformat{\subsection}{\large\bfseries\color{fgsoft}}{\thesubsection}{0.6em}{}
\titlespacing*{\subsection}{0pt}{1.3em}{0.4em}
\titleformat{\subsubsection}[runin]{\bfseries\color{fgsoft}}{}{0pt}{}[.]

\pagestyle{fancy}
\fancyhf{}
\lhead{\footnotesize\color{muted}Manual de usuario}
\rhead{\footnotesize\color{muted}\thepage}
\renewcommand{\headrulewidth}{0.4pt}
\renewcommand{\headrule}{\color{line}\hrule height 0.4pt}

\setlength{\parindent}{0pt}
\setlength{\parskip}{0.62em}

% Enlaces y referencias en el mismo gris del cuerpo: sin subrayados ni azules.
\hypersetup{colorlinks=false}

% ── Recuadros ───────────────────────────────────────────────────────────────
% Se distinguen por una barra lateral y el peso del titulo, no por color.
\newtcolorbox{nota}[1][Nota]{
  enhanced, breakable, colback=surfacealt, colframe=surfacealt,
  borderline west={2pt}{0pt}{line},
  coltitle=fg, fonttitle=\bfseries\footnotesize\sffamily,
  title=#1, boxrule=0pt, arc=0pt, outer arc=0pt,
  left=10pt, right=10pt, top=7pt, bottom=7pt, toptitle=1pt, bottomtitle=3pt}

\newtcolorbox{aviso}[1][Importante]{
  enhanced, breakable, colback=surface, colframe=surface,
  borderline west={2pt}{0pt}{fg},
  coltitle=fg, fonttitle=\bfseries\footnotesize\sffamily,
  title=#1, boxrule=0pt, arc=0pt, outer arc=0pt,
  left=10pt, right=10pt, top=7pt, bottom=7pt, toptitle=1pt, bottomtitle=3pt}

% ── Captura de pantalla con marcador de posicion ────────────────────────────
% Uso: \captura{nombre-archivo}{Pie de la figura}
% Si manual/capturas/<nombre>.png existe, la inserta. Si no, dibuja el hueco.
\newcommand{\captura}[2]{%
  \begin{figure}[htbp]
  \centering
  \IfFileExists{capturas/#1.png}{%
    % Ancho completo, pero acotado en alto: una captura de ventana completa
    % puede ser muy alta y desbordaria la pagina. keepaspectratio evita que se
    % deforme cuando manda la restriccion de altura.
    {\color{line}\fboxrule=0.6pt\fboxsep=0pt\fbox{%
      \includegraphics[width=0.94\textwidth, height=0.62\textheight, keepaspectratio]{capturas/#1.png}}}%
  }{%
    \fcolorbox{line}{surfacealt}{%
      \begin{minipage}[c][3.1cm][c]{0.93\textwidth}
        \centering
        {\footnotesize\bfseries\color{muted}CAPTURA PENDIENTE}\\[5pt]
        {\ttfamily\footnotesize\color{fgsoft} capturas/#1.png}\\[4pt]
        {\footnotesize\color{muted} #2}
      \end{minipage}}%
  }
  \caption{#2}
  \label{fig:#1}
  \end{figure}%
}

% Pies de figura discretos
\usepackage{caption}
\captionsetup{font={small,color=muted}, labelfont={bf,color=muted}, skip=6pt}

 desde la
% carpeta de cada manual.
%
% La paleta sale de las variables de frontend/src/style.css: la interfaz es
% enteramente acromatica, y los manuales tampoco introducen color.
% ================================================================================
\documentclass[11pt, letterpaper]{article}

\usepackage[spanish,mexico]{babel}
\usepackage[utf8]{inputenc}
\usepackage[T1]{fontenc}

% Tipografia: Inter, la mas cercana en TeX Live a la Geist que usa la
% aplicacion. Sin serifas en todo el documento, igual que la interfaz.
\usepackage[sfdefault]{inter}
\usepackage[scaled=0.86]{beramono}

\usepackage[top=2.4cm, bottom=2.4cm, left=2.6cm, right=2.6cm]{geometry}
\usepackage{graphicx}
\usepackage{xcolor}
\usepackage{enumitem}
\usepackage{fancyhdr}
\usepackage{titlesec}
\usepackage{tcolorbox}
% `enhanced` y `borderline west` viven en skins; `breakable` en su propia
% libreria. Sin declararlas, tcolorbox ignora esas claves en silencio y los
% recuadros pierden la barra lateral que los distingue.
\tcbuselibrary{skins,breakable}
\usepackage{booktabs}
\usepackage[hidelinks]{hyperref}

% ── Paleta ──────────────────────────────────────────────────────────────────
% Tomada de las variables de src/style.css. La interfaz es enteramente
% acromatica —todos sus colores son oklch(L 0 0)— asi que el manual tampoco
% introduce color: solo grises, con el texto como unico elemento de contraste.
\definecolor{fg}{HTML}{0A0A0A}        % --foreground
\definecolor{fgsoft}{HTML}{171717}    % --primary
\definecolor{muted}{HTML}{737373}     % --muted-foreground
\definecolor{surface}{HTML}{F5F5F5}   % --muted / --secondary
\definecolor{surfacealt}{HTML}{FAFAFA}% --sidebar
\definecolor{line}{HTML}{E5E5E5}      % --border

\titleformat{\section}{\LARGE\bfseries\color{fg}}{\thesection}{0.7em}{}
\titlespacing*{\section}{0pt}{1.6em}{0.7em}
\titleformat{\subsection}{\large\bfseries\color{fgsoft}}{\thesubsection}{0.6em}{}
\titlespacing*{\subsection}{0pt}{1.3em}{0.4em}
\titleformat{\subsubsection}[runin]{\bfseries\color{fgsoft}}{}{0pt}{}[.]

\pagestyle{fancy}
\fancyhf{}
\lhead{\footnotesize\color{muted}Manual de usuario}
\rhead{\footnotesize\color{muted}\thepage}
\renewcommand{\headrulewidth}{0.4pt}
\renewcommand{\headrule}{\color{line}\hrule height 0.4pt}

\setlength{\parindent}{0pt}
\setlength{\parskip}{0.62em}

% Enlaces y referencias en el mismo gris del cuerpo: sin subrayados ni azules.
\hypersetup{colorlinks=false}

% ── Recuadros ───────────────────────────────────────────────────────────────
% Se distinguen por una barra lateral y el peso del titulo, no por color.
\newtcolorbox{nota}[1][Nota]{
  enhanced, breakable, colback=surfacealt, colframe=surfacealt,
  borderline west={2pt}{0pt}{line},
  coltitle=fg, fonttitle=\bfseries\footnotesize\sffamily,
  title=#1, boxrule=0pt, arc=0pt, outer arc=0pt,
  left=10pt, right=10pt, top=7pt, bottom=7pt, toptitle=1pt, bottomtitle=3pt}

\newtcolorbox{aviso}[1][Importante]{
  enhanced, breakable, colback=surface, colframe=surface,
  borderline west={2pt}{0pt}{fg},
  coltitle=fg, fonttitle=\bfseries\footnotesize\sffamily,
  title=#1, boxrule=0pt, arc=0pt, outer arc=0pt,
  left=10pt, right=10pt, top=7pt, bottom=7pt, toptitle=1pt, bottomtitle=3pt}

% ── Captura de pantalla con marcador de posicion ────────────────────────────
% Uso: \captura{nombre-archivo}{Pie de la figura}
% Si manual/capturas/<nombre>.png existe, la inserta. Si no, dibuja el hueco.
\newcommand{\captura}[2]{%
  \begin{figure}[htbp]
  \centering
  \IfFileExists{capturas/#1.png}{%
    % Ancho completo, pero acotado en alto: una captura de ventana completa
    % puede ser muy alta y desbordaria la pagina. keepaspectratio evita que se
    % deforme cuando manda la restriccion de altura.
    {\color{line}\fboxrule=0.6pt\fboxsep=0pt\fbox{%
      \includegraphics[width=0.94\textwidth, height=0.62\textheight, keepaspectratio]{capturas/#1.png}}}%
  }{%
    \fcolorbox{line}{surfacealt}{%
      \begin{minipage}[c][3.1cm][c]{0.93\textwidth}
        \centering
        {\footnotesize\bfseries\color{muted}CAPTURA PENDIENTE}\\[5pt]
        {\ttfamily\footnotesize\color{fgsoft} capturas/#1.png}\\[4pt]
        {\footnotesize\color{muted} #2}
      \end{minipage}}%
  }
  \caption{#2}
  \label{fig:#1}
  \end{figure}%
}

% Pies de figura discretos
\usepackage{caption}
\captionsetup{font={small,color=muted}, labelfont={bf,color=muted}, skip=6pt}



\begin{document}

% ── Portada ─────────────────────────────────────────────────────────────────
\begin{titlepage}
\thispagestyle{empty}
\vspace*{5.5cm}
{\fontsize{46}{50}\selectfont\bfseries\color{fg} GAMMA\par}
\vspace{0.5em}
{\large\color{muted} Gobierno Automatizado del Maestro de Materiales\par}

\vspace{2.6cm}
{\color{line}\hrule height 1pt}
\vspace{1.2em}

{\huge\bfseries\color{fg} Manual de instalación\par}
\vspace{0.5em}
{\large\color{muted} Puesta en marcha, actualización y operación\par}

\vfill
{\color{line}\hrule height 0.4pt}
\vspace{1em}
{\footnotesize\color{muted}
Unidad de Energía \quad·\quad \today\par
\vspace{0.4em}
Dirigido a quien despliega y mantiene la plataforma.
El uso diario se documenta en el manual de usuario, y la administración en el
manual técnico.\par}
\vspace{1cm}
\end{titlepage}

\tableofcontents
\newpage

% ============================================================================
\section{Requisitos}

\begin{itemize}[leftmargin=*, itemsep=3pt]
  \item \textbf{Docker} con el complemento Compose. Es el único requisito
        obligatorio: todo lo demás corre dentro de contenedores.
  \item \textbf{Git}, para obtener y actualizar el código.
  \item Una \textbf{clave de acceso al servicio de modelo de lenguaje}. Es el
        único componente que no es de código abierto y el único que tiene costo,
        bajo un esquema de pago por uso.
  \item Alrededor de \textbf{4 GB de memoria disponible}. El modelo de
        clasificación se carga completo en memoria, y durante un reentrenamiento
        conviven dos en paralelo.
\end{itemize}

\newpage
% ============================================================================
\section{Servicios}

Todo se orquesta con Docker Compose sobre una red interna.

\begin{center}
\begin{tabular}{@{}llp{6.5cm}@{}}
\toprule
\textbf{Servicio} & \textbf{Puerto} & \textbf{Qué es} \\
\midrule
\texttt{db}       & 5436 & PostgreSQL 16 con la extensión \texttt{pg\_trgm}, que resuelve la búsqueda por similitud \\
\texttt{api}      & 8001 & FastAPI sobre Python 3.12. Incluye el modelo de clasificación y la conexión al modelo de lenguaje \\
\texttt{lab}      & 8888 & Jupyter Lab con los cuadernos de análisis \\
\texttt{frontend} & 8089 & La aplicación web. Sirve además de proxy hacia los demás servicios \\
\bottomrule
\end{tabular}
\end{center}

\subsection{El frontend como punto de entrada}

En producción, el frontend no solo sirve la aplicación: unifica el acceso a todo
bajo un mismo dominio.

\begin{center}
\begin{tabular}{@{}ll@{}}
\toprule
\textbf{Ruta} & \textbf{Destino} \\
\midrule
\texttt{/}           & La aplicación web \\
\texttt{/api/*}      & El API \\
\texttt{/swagger-ui} & La documentación interactiva del API \\
\texttt{/lab/*}      & Jupyter Lab \\
\bottomrule
\end{tabular}
\end{center}

\subsection{Los dos modos del frontend}

El comportamiento lo decide la variable \texttt{FRONTEND\_MODE}:

\begin{itemize}[leftmargin=*, itemsep=3pt]
  \item \textbf{development} — servidor de desarrollo con recarga automática.
        Los cambios en el código se reflejan sin reconstruir nada.
  \item \textbf{production} — compila la aplicación y sirve archivos estáticos.
        Más rápido, pero \textbf{cada cambio exige reconstruir la imagen}.
\end{itemize}

\begin{nota}[Cuál usar]
En modo desarrollo una actualización se aplica con solo recrear el contenedor,
lo que simplifica bastante el mantenimiento. Si opta por modo producción,
recuerde que tras cada actualización hay que reconstruir:
\texttt{docker compose up -d -{}-build frontend}.
\end{nota}

\newpage
% ============================================================================
\section{Variables de entorno}

Se definen en un archivo \texttt{.env} en la raíz del proyecto. Parta de
\texttt{.env.example}.

\begin{center}
\begin{tabular}{@{}lp{8.4cm}@{}}
\toprule
\textbf{Variable} & \textbf{Para qué sirve} \\
\midrule
\texttt{DB\_USER}, \texttt{DB\_PASSWORD} & Credenciales de la base de datos \\
\texttt{DB\_NAME}, \texttt{DB\_HOST}, \texttt{DB\_PORT} & Ubicación de la base \\
\texttt{JWT\_SECRET} & Clave con la que se firman las sesiones. \textbf{Debe ser un valor largo y aleatorio} \\
\texttt{JWT\_EXPIRY\_HOURS} & Duración de la sesión \\
\texttt{LLM\_API\_KEY} & Clave del servicio de modelo de lenguaje \\
\texttt{LLM\_MODEL} & Modelo de lenguaje a utilizar \\
\texttt{JUPYTER\_TOKEN} & Token de acceso al laboratorio \\
\texttt{CONFIANZA\_UMBRAL} & A partir de qué confianza una sugerencia se da por resuelta \\
\texttt{DUPLICADO\_UMBRAL} & A partir de qué similitud se alerta de un posible duplicado \\
\texttt{FRONTEND\_MODE} & \texttt{development} o \texttt{production} \\
\bottomrule
\end{tabular}
\end{center}

\begin{aviso}[No versione el archivo .env]
Contiene credenciales y claves. Está excluido del repositorio a propósito. Cada
entorno mantiene el suyo.
\end{aviso}

\subsection{Los dos umbrales}

Merecen atención porque cambian el comportamiento del asistente:

\begin{itemize}[leftmargin=*, itemsep=3pt]
  \item \texttt{DUPLICADO\_UMBRAL} está deliberadamente bajo. Bajarlo aún más
        genera más alertas irrelevantes; subirlo deja pasar duplicados reales,
        que son mucho más costosos.
  \item \texttt{CONFIANZA\_UMBRAL} decide cuándo el sistema da por buena su
        propia sugerencia. Subirlo aumenta la exactitud pero deriva más casos a
        revisión manual.
\end{itemize}

\newpage
% ============================================================================
\section{Primera instalación}

\subsection{Obtener el código y configurar}

\begin{verbatim}
git clone <repositorio> gamma
cd gamma
cp .env.example .env
\end{verbatim}

Complete el \texttt{.env} con las credenciales y claves.

\subsection{Levantar los servicios}

\begin{verbatim}
docker compose up -d --build
\end{verbatim}

La primera vez tarda varios minutos: descarga las imágenes base e instala
dependencias.

\subsection{Aplicar las migraciones}

\begin{verbatim}
docker compose exec api python migrate.py --status
docker compose exec api python migrate.py
\end{verbatim}

El primer comando lista qué hay aplicado y qué falta; el segundo aplica lo
pendiente. Es idempotente: repetirlo no causa daño.

\begin{nota}[Los scripts de inicialización solo corren una vez]
Los archivos de \texttt{db/init/} se ejecutan únicamente cuando el volumen de
datos está vacío. En una base ya creada no vuelven a correr: para eso están las
migraciones.
\end{nota}

\subsection{Crear el primer administrador}

Registre una cuenta desde la aplicación y promuévala. \textbf{Este paso es
obligatorio}: sin él nadie puede entrar a las secciones de administración.

\begin{verbatim}
docker compose exec db psql -U <usuario> -d <base> \
  -c "UPDATE public.users SET admin=true WHERE email='...';"
\end{verbatim}

A partir de ahí, los demás roles se gestionan desde la sección Usuarios.

\subsection{Cargar los datos iniciales}

Desde la sección \textbf{Datos}, pestaña Importar, cargue en este orden: UNSPSC,
luego el catálogo de clases y al final los archivos del maestro. El orden
importa; el manual técnico explica por qué.

\newpage
% ============================================================================
\section{Actualizar una instalación existente}

\subsection{Procedimiento}

\begin{verbatim}
cd /ruta/del/proyecto
git status          # revisar que no haya cambios locales sin guardar
git pull
docker compose exec api python migrate.py
docker compose up -d --force-recreate --no-deps frontend
\end{verbatim}

\subsection{Por qué cada paso}

\begin{itemize}[leftmargin=*, itemsep=4pt]
  \item \textbf{Las migraciones} deben correr siempre. Una pantalla que consulta
        una tabla que todavía no existe falla sin explicación clara.
  \item \textbf{Recrear el frontend} es necesario porque algunos archivos se
        montan individualmente en el contenedor, y al reemplazarlos el
        contenedor pierde la referencia. Se detalla más adelante.
  \item \textbf{El API no requiere nada}: su código se monta como directorio y
        el servidor recarga solo. Solo hay que reconstruirlo si cambiaron sus
        dependencias.
\end{itemize}

\begin{aviso}[Reconstruir si cambian las dependencias]
Si la actualización modificó las dependencias del API o de la aplicación web, sí
hay que reconstruir la imagen correspondiente con
\texttt{docker compose up -d -{}-build}. Revise el registro de cambios ante la
duda.
\end{aviso}

\newpage
% ============================================================================
\section{Problemas frecuentes}

\subsection*{Tras actualizar, la aplicación muestra un error de archivo no encontrado}

Es el problema más común y no significa que falte nada. Algunos archivos se
montan individualmente en el contenedor, y Docker los sigue por su identificador
interno, no por su nombre. Cuando la actualización reemplaza uno de esos
archivos, el contenedor queda apuntando a algo que ya no existe.

\begin{verbatim}
docker compose up -d --force-recreate --no-deps frontend
\end{verbatim}

\begin{nota}[Reiniciar no basta]
\texttt{docker compose restart} reutiliza el mismo contenedor y conserva el
montaje defectuoso. Hay que recrearlo.
\end{nota}

\subsection*{El pull falla diciendo que las ramas divergieron}

Ocurre si el historial del repositorio se reescribió. En un servidor, que solo
consume código y no produce cambios, la salida es adoptar el remoto tal cual:

\begin{verbatim}
git fetch origin
git reset --hard origin/main
\end{verbatim}

\begin{aviso}[Revise antes de descartar]
\texttt{git reset -{}-hard} descarta cualquier cambio local en archivos
versionados. Ejecute \texttt{git status} primero. Los archivos excluidos del
repositorio, como \texttt{.env} y los artefactos del modelo, no se ven
afectados.
\end{aviso}

\subsection*{El pull falla por archivos sin versionar que serían sobrescritos}

Sucede cuando un archivo que antes estaba excluido pasa a formar parte del
repositorio. Aparta los locales y vuelve a intentar:

\begin{verbatim}
mkdir -p ~/respaldo && mv <archivos> ~/respaldo/
git pull
\end{verbatim}

\subsection*{Nadie puede entrar a las secciones de administración}

Ninguna cuenta tiene el rol. Promueva una desde la base de datos, como se
describe en la primera instalación.

\subsection*{Una pantalla falla al consultar datos}

Casi siempre son migraciones sin aplicar. Verifique con
\texttt{migrate.py -{}-status}.

\subsection*{El asistente no responde}

Revise que \texttt{LLM\_API\_KEY} esté configurada y que el servicio externo
esté disponible. Los errores de comunicación quedan registrados y se pueden
consultar desde la sección Datos.

\newpage
% ============================================================================
\section{Respaldos}

Tres cosas merecen respaldo, y solo una está en el repositorio:

\begin{center}
\begin{tabular}{@{}lp{8.6cm}@{}}
\toprule
\textbf{Qué} & \textbf{Por qué} \\
\midrule
La base de datos & Contiene el catálogo, las solicitudes, las conversaciones y las cuentas. Es lo único verdaderamente irrecuperable \\
El archivo \texttt{.env} & Credenciales y claves. No está versionado \\
Los artefactos del modelo & No viajan por el repositorio. Se pueden regenerar reentrenando, pero eso toma tiempo y no reproduce exactamente el modelo anterior \\
\bottomrule
\end{tabular}
\end{center}

El código no necesita respaldo: vive en el repositorio.

\begin{nota}[Los artefactos son propios de cada entorno]
Como no se versionan, cada instalación mantiene los suyos. Un entorno recién
desplegado arranca con el modelo inicial y su propio historial de versiones.
\end{nota}

\vspace{1cm}

\begin{aviso}[Antes de recrear la base]
\texttt{docker compose down -v} borra los volúmenes, y con ellos toda la
información. Úselo solo cuando quiera partir de cero, y nunca sin respaldo
previo.
\end{aviso}

\end{document}
